\begin{exercisechunk}
\item \textbf{When does a reusable jump table reduce the total cost?}\suppmark
\label[exercise]{ex:l14-pointer-doubling}
\exerciseprofile{2}{2}{2}{\exproof\quad\excalculation\quad\exinterpretation}{%
Count the preprocessing and query operations that the asymptotic depth
comparison suppresses.}
Let $p:\cD\to\cD$ be a self-map on a set of size $M$. Suppose full-attention
preprocessing has constructed
\[
 q_\ell(u)=p^{2^\ell}(u),
 \qquad \ell=0,\ldots,J,
\]
at every map-record position. Constructing each of
$q_1,\ldots,q_J$ costs $M$ table updates.
Count one application of a stored table entry to one key as one table lookup,
and count the construction of one new $q_\ell(u)$ value as one table update.
\begin{enumerate}[(a),beginpenalty=10000]
\item Let $K=\sum_{\ell=0}^J b_\ell2^\ell$, where
$b_\ell\in\{0,1\}$. Prove that applying the maps $q_\ell$ for which
$b_\ell=1$ returns $p^K(x)$ in any order. Deduce that the query uses
$\sum_{\ell=0}^J b_\ell$ lookups.
\item Consider the exhaustive batch containing one query for every
$K\in\{0,\ldots,2^{J+1}-1\}$. Compute exactly the total number of lookups
used by direct iteration and by the stored jump tables. Your answer should
include the $MJ$ preprocessing updates in the jump-table total.
\item Suppose instead that the same input map must answer $Q$ queries, each
with iteration count
$K_{\max}=2^{J+1}-1$ but possibly with a different initial state. Compute
the total jump-table cost and the total direct-iteration cost. State an
inequality on $M,Q,K_{\max}$ under which the jump tables use fewer total
operations.
\item Compare the exhaustive batch in part~(b) with the maximum-length batch
in part~(c). Explain why a sequential-depth advantage
does not by itself imply a reduction in the total number of table
operations.
\end{enumerate}
\end{exercisechunk}
